\tjformsection{一、毕业论文的课题背景}

视觉语言导航（Vision-and-Language Navigation，VLN）要求智能体根据自然语言指令，在持续感知周围环境的过程中完成自主导航，是具身智能连接语言理解、环境感知与机器人决策的重要任务。随着视觉语言大模型和视频大模型的发展，基于连续视觉观测直接生成导航动作成为新的研究方向。然而，现有方法大多以RGB视觉为主要感知来源，在真实机器人环境中容易受到遮挡、光照变化、视角变化及空间尺度不确定等因素影响；同时，长时间连续输入还会带来历史信息不断增长、重复计算和推理延迟增加等问题。

近年来，流式视觉语言导航开始关注连续观测条件下的低延迟决策与长期上下文建模。StreamVLN等工作通过滑动窗口维护近期视觉上下文，并利用慢更新记忆压缩历史视觉信息，使模型能够在有限上下文规模下处理长时间视频流；其进一步利用深度和里程计将视觉特征反投影到三维空间，实现无需额外训练的空间冗余裁剪。该类方法为流式导航提供了良好基础，但深度、几何结构和机器人运动状态仍主要作为辅助信息使用，尚未充分形成统一的多模态在线感知与时空记忆机制。

本课题拟面向连续机器人导航场景，研究多模态流式感知与时空记忆驱动的视觉语言导航方法。重点探索在保持已有视觉语言或导航模型主体参数固定的条件下，融合RGB视觉语义、深度/点云几何以及由IMU、视觉里程计或激光里程计获得的机器人运动状态，构建严格因果的在线输入、轻量级多模态融合和有界长期记忆机制。在避免大规模端到端训练的前提下，提高模型对长时导航、复杂环境及感知退化条件的适应能力，并通过公开连续VLN基准进行实验验证。

\tjformsection{二、毕业论文应完成的主要内容}

\begin{enumerate}
  \item 调研视觉语言导航、连续环境导航、视觉语言大模型、多模态机器人感知、流式视频理解、长时记忆与Cache等相关研究，梳理现有方法的任务设定、技术路线和主要局限，明确课题的研究问题、实验范围以及“冻结大模型主体、轻量级扩展”的总体技术路线。
  \item 研究R2R-CE、RxR-CE等连续视觉语言导航数据集及Habitat仿真平台的数据组织、导航指令、传感器观测和评测方式，建立严格因果的流式数据处理与在线评测流程。根据实验条件组织RGB、Depth、位姿/里程计等多源观测，并可由深度信息生成局部点云或模拟激光雷达几何观测，为后续多模态研究提供统一输入。
  \item 选取具有代表性的视觉语言导航或视频视觉语言模型作为基础模型，复现或调用现有模型建立视觉流式导航基线。重点分析单帧、固定历史窗口和流式历史输入等不同设置，建立Success Rate、SPL、Navigation Error等统一评测流程，并保证推理过程中不访问未来观测。
  \item 研究面向流式导航的多模态感知表示与轻量级融合方法。在RGB视觉语义基础上，引入深度/点云描述的局部三维几何、可通行空间以及机器人位姿和运动状态，探索结构化几何描述、特征投影、小型Adapter或其他参数高效方法，将多源信息引入导航决策，并比较“仅视觉”“视觉+几何”“视觉+几何+运动状态”等方案。
  \item 研究适用于长时连续导航的时空记忆机制。针对视觉流持续到达时历史信息冗余、上下文长度增长及空间重复观测等问题，探索滑动窗口、状态缓存、Memory Bank、特征Cache、KV Cache、基于三维位置的空间裁剪或语义—几何记忆等方式，在保持记忆规模有界的条件下利用近期和长期历史信息，提高长时导航中的空间一致性与决策稳定性。
  \item 在R2R-CE、RxR-CE等数据集上开展对比实验与消融实验，从导航成功率、路径效率、长时记忆能力、在线推理延迟、显存/Token开销以及多模态信息贡献等方面评价方法。
\end{enumerate}

\tjformsection{三、毕业论文的基本要求及应完成的成果形式}

\tjformsubsection{基本要求}

\begin{enumerate}
  \item 掌握视觉语言导航、视觉语言模型、连续环境导航、多模态感知、时序建模与流式推理等基本原理，具备相关文献检索、阅读和分析能力。
  \item 熟悉Habitat及R2R-CE、RxR-CE等连续导航数据，能够完成语言指令、RGB/Depth观测、位姿信息及导航轨迹的数据预处理、同步组织和在线输入构建。
  \item 能够复现或使用具有代表性的VLN/VLM基础模型，建立统一的视觉流式导航实验流程；原则上保持大模型主体参数固定，优先采用training-free、轻量级Adapter或少量参数更新的方法，避免依赖大规模端到端训练。
  \item 完成至少一种多模态流式输入融合方案和一种时空记忆/历史信息管理方案的设计、实现与验证，并通过对比实验和消融实验分析视觉、几何、运动状态及记忆模块的实际作用。
  \item 实验过程应记录完整，实验指标、运行环境和参数配置应清晰可复现；对实验结果进行客观分析，不仅关注导航成功率与路径效率，也应评价模型的在线推理延迟、显存/Token开销以及复杂环境下的鲁棒性。
  \item 按照学校要求完成毕业论文及相关材料，准确说明研究问题、方法设计、实验过程及主要结论。
\end{enumerate}

\tjformsubsection{成果形式}

\begin{enumerate}
  \item 毕业论文一份。
  \item VLN数据处理、多模态输入组织、流式推理、时空记忆及实验评测相关程序代码。
  \item 模型配置文件、运行说明及实验复现文档。
  \item 开题报告、中期检查材料、毕业答辩PPT及其他学校要求的相关材料。
\end{enumerate}

\tjformsection{四、毕业论文的进度安排}

\begin{enumerate}
  \item 第1—4周：开展视觉语言导航、连续环境导航、视觉语言大模型、多模态感知及流式记忆相关文献调研，熟悉R2R-CE、RxR-CE和Habitat等数据与平台，明确研究问题和实验设定，完成开题准备。
  \item 第5—8周：搭建实验环境，完成数据预处理和传感器观测读取流程，复现或运行代表性的VLN/VLM模型，建立基于视觉输入的导航基线和统一评测脚本。
  \item 第9—12周：将基础模型整理为严格因果的流式推理模式，构建单帧、固定窗口和历史流输入等baseline，分析历史长度、更新频率以及缓存策略对导航精度和推理效率的影响。
  \item 第13—18周：研究多模态流式感知与轻量级融合方案，引入Depth/局部点云、可通行空间及机器人位姿/运动状态，完成“视觉+几何”“视觉+几何+状态”等主要方案并开展初步实验。
  \item 第19—24周：研究短时上下文与长期时空记忆机制，探索滑动窗口、特征Cache、KV Cache、基于三维位置的空间裁剪或语义—几何记忆等方法，完成主要方法设计及核心实验。
  \item 第25—27周：补充不同数据集、不同历史长度和不同模态组合下的对比实验与消融实验，分析感知退化、几何噪声或里程计扰动等条件下的鲁棒性，并根据条件探索真实机器人扩展。
  \item 第28—30周：整理实验结果、方法结构图及分析图表，完成毕业论文初稿，并整理代码、配置文件和复现实验说明。
  \item 第31—32周：根据指导意见修改完善毕业论文，完成最终材料提交和毕业答辩准备。
\end{enumerate}

\tjformsection{五、毕业论文应收集的资料及主要参考文献}

应重点收集以下资料：视觉语言导航与连续环境导航相关论文及开源代码；R2R-CE、RxR-CE、Matterport3D、HM3D及Habitat平台资料；Video-LLM/VLA在机器人导航中的应用研究；RGB-D、点云、里程计与机器人状态等多模态感知方法；Memory Bank、KV Cache、Token Pruning、空间记忆及流式推理相关方法；VLN评价指标、实验协议以及GPU推理延迟、显存和Token开销分析工具等。

主要参考文献建议包括：

% 编号顺序即 \nocite 的顺序（GB/T 7714 顺序编码制）
\nocite{anderson2018vln,krantz2020vlnce,zhang2024navid,zhang2025uninavid,%
        cheng2025navila,wei2025streamvln,zhang2025mapnav,krantz2021waypoint,%
        ramakrishnan2021hm3d}
\printbibliography[heading=none]

\tjformsection{六、其他要求}

定期向指导教师汇报进展，根据调研和实验结果完善研究计划；优先保证流式推理的因果性、实验可复现性和方法的轻量化，保留程序代码、实验配置及原始结果，按时完成各阶段任务。

\vfill
% 任务书正文末尾的签字行为四号黑体（审核意见框内的签字行为小四号，见开题/中期报告）
\tjformsignature[\zihao{4}]{指导教师签名}{9em}
